\subsection*{The Partial Type: \textbf{mdspan}}

An \texttt{mdspan} is a partial type: it describes a shape but cannot by itself declare a data buffer. Shape values can nevertheless be defined and manipulated independently. Tiling depends on multidimensional shape rather than element type, so this separation simplifies shape calculations.

The compiler also checks \texttt{mdspan} expressions to detect errors as early as possible, before execution when the necessary information is available.

In Croqtile, a mdspan is braced by '[' and ']'. The upper-bound of dimensions inside are separated by ','. That is,

\begin{lstlisting}[language=croqtile,caption={An mdspan declaration}]
mdspan sp : [7, 8];           // defined a mdspan of 2-dimensions
mdspan<4> mds : [a, 3, 4, 1]; // 'a' is an existing integer
d : [c, 4, 28];               // 'd' is not explicit annotated. Type is deduced.
\end{lstlisting}

An \texttt{mdspan} can be declared explicitly, with an optional rank in angle brackets, or inferred from its initializer, as in the declaration of \texttt{d}. The types of \texttt{mds} and \texttt{d} depend on the values of \texttt{a} and \texttt{b}. If these values are supplied at runtime, some checks must also occur at runtime.

For a C++ programmer, an \texttt{mdspan} describes the ranges of a multidimensional array. In the example, \texttt{sp} has two dimensions, with indices from 0 to 6 and from 0 to 7, respectively.

When a mdspan is defined, it is possible to get the integer dimension value using operator '()' over the mdspan.

\begin{lstlisting}[language=croqtile,caption={A shape literal}]
sp : [7, 8];
int b = sp(0) + sp(1);   // 'b' equals to 15 (7 + 8)
\end{lstlisting}

The expression \co@sp(0)@ returns the first extent of \co@sp@. Extents are integer values and can participate in integer arithmetic.

You can therefore derive a new \texttt{mdspan} from an existing one:

\begin{lstlisting}[language=croqtile,caption={A shape with one dimension changed}]
sp : [6, 8];
spn : [1, sp(0)/2, sp(1)/4]; // define a new mdspan from the existing one.
\end{lstlisting}

This makes tiling multidimensional shapes convenient. Because tiling is common in high-performance kernels, Croqtile provides a shorter form for the same calculation:

\begin{lstlisting}[language=croqtile,caption={The same shape again}]
sp : [6, 8];
spn : sp [1, (0)/2, (1)/4];  // spn is defined as [1, 3, 2]
\end{lstlisting}

This version computes the same shape with less code, following Croqtile's aim of expressing common operations concisely.

A named \co@mdspan@ must be initialized when declared and cannot be modified afterwards. To describe a different shape, derive a new name from the existing one. This preserves the relationship between shape expressions for the compiler's analysis.

An \texttt{mdspan} can also be defined using operations on whole tuples. The \texttt{ituple} discussion explains those forms.
